\section*{Chicken Tikka Masala}
\ihead{}\chead{Personen:2}\ohead{}
\ifoot{Vorbereitungszeit:20 min}\cfoot{}\ofoot{Kochzeit:20 min}
{\Large Zutaten}
\begin{multicols}{2}
\begin{itemize}
\item \SI{400}{g} Hähnchenbrustfilet
\item \SI{100}{g} Naturjoghurt
\item \num{1} kleine Dose Tomatenmark, konzentriert
\item \num{2}{TL} Ingwer, gerieben
\item \num{1}{TL} Knoblauch, gepresst (nach Belieben mehr)
\item \num{1}{TL} Currypulver
\item \num{1}{TL} Paprikapulver, edelsüß
\item \num{2}{Msp.} Korianderpulver
\item \num{2}{Msp.} Kreuzkümmelpulver
\item \num{1}{Msp.} Nelkenpulver
\item \num{1}{Msp.} Zimtpulver
\item \num{1}{Msp.} Cayennepfeffer, nach Belieben mehr
\item \num{1} Prise Vanillepulver
\item etwas Öl
\item Salz
\item \num{1} große Schalotte oder \num{1} kleine Zwiebel
\item Paprikaschote nach Bedarf
\item Zitronensaft nach Bedarf
\end{itemize}
% \columnbreak
\end{multicols}
\noindent {\Large Zubereitung}\ \
\textit{Vorbereitung:} Die Hühnerbrust in mundgerechte Stückchen schneiden und in eine verschließbare Schüssel oder einen ausreichend großen Gefrierbeutel geben.
Joghurt, Tomatenmark und alle Gewürze – also Knoblauch, Ingwer, Zimt, Nelken, Curry, Paprika, Koriander, Kreuzkümmel, Vanillepulver und Cayennepfeffer – zugeben.
Alles gut vermengen, sodass das Fleisch gut von der Marinade überzogen ist.
Die Schüssel bzw. den Gefrierbeutel verschließen und für mehrere Stunden, eventuell über Nacht, in den Kühlschrank geben, sodass die Gewürze schön in das Fleisch einziehen können.
Ist das Fleisch gut durchgezogen, die Schalotte abziehen, in feine Würfel schneiden, salzen und in wenig Öl goldgelb anbraten.
Durch das Salz wird die Schalotte weich, zerfällt später in der Sauce und gibt eine schöne Bindung.
Nun das Fleisch samt Marinade zugeben und bei mittlerer Hitze ca. \SI{15}{min} garen, bis das Hühnerfleisch durchgebraten ist.
Das Fleisch bleibt durch den Joghurt schön saftig.
Sollte die Sauce zu dick oder zu wenig sein, einfach etwas Wasser dazugeben.
Nach dem Abschmecken mit körnigem Reis oder Naan-Brot und einem frischen Salat servieren.
Nach Belieben kann eine gewürfelte Paprikaschote hinzugegeben werden.
Ein Spritzer Zitronensaft verleiht dem Gericht eine feinsäuerliche, fruchtige Note.
